Denote by $S_0$ the reduced subscheme of $S$ whose underlying space is the closed point of $S$. Let $T'$ be an element of $\Ker[M(S')\rightrightarrows M(S'')]$, $T''$ its image in $M(S'')$. Since $S'_0=S'\times_S S_0$ is faithfully flat over $S_0$, there exists an $S$-subgroup of multiplicative type $T_0$ of $G_{S_0}$ whose inverse image in $G_{S'_0}$ is $T'_{S'_0}$ (fpqc descent). But $S$ is local noetherian complete, so there exist an $S$-group $T$ of multiplicative type and an $S_0$-morphism $u_0:T\times_S S_0\to T_0$ (Exp. X 3.3). The inverse image $u'_0$ of $u_0$ over $S'_0$ extends uniquely to an $S'$-isomorphism $u':T_{S'}\to T'$, again by Exp. X 3.3 (note that $S'$, being finite over $S$ local complete, is the sum of finitely many complete local schemes). The two inverse images of $u'$ over $S''$ are two morphisms from $T_{S''}$ to $T''$ which agree on $S_0\times_S S''$, hence agree (\loccit). Since $T$ is flat over $S$ and $S'\to S$ is a finite effective epimorphism, it follows from TDTE I page 8 that the diagram:
\[
\Hom_S(T,G)\to\Hom_{S'}(T_{S'},G_{S'})\rightrightarrows\Hom_{S''}(T_{S''},G_{S''})
\]
is exact. Thus $u'$ comes from a morphism of preschemes $u:T\to G$. Likewise, $T\times_S T$ is flat over $S$, and consequently the map:
\[
\Hom_S(T\times_S T,G)\to\Hom_{S'}(T_{S'}\times_{S'}T_{S'},G_{S'})
\]
is injective. One immediately deduces that $u$ is a group morphism, since this is so for $u'$. Moreover, $u$ is a monomorphism. Indeed, note first that $\Ker u$ is flat over $S$, for to establish this fact one may suppose $S$ local artinian (EGA $0_{\mathrm{III}}$ 10.2.2), hence $G$ separated (Exp. VIA 0.3), but then $\Ker u$ is of multiplicative type (Exp. IX 6.8), hence is flat over $S$. Since $\Ker(u)\times_S S'=\Ker(u')=0$, one has $\Ker(u)=0$ (3.8). But $u$, being a monomorphism, is an immersion (Exp. VIII, Remarks 7.13), and the image group $u(T)$ is indeed an element of $M(S)$ whose image in $M(S')$ is $T'$, which completes the proof of 3.7.

\noindent\textbf{e)} End of the proof of I).

By reduction c), one is reduced to the case where $S$ is the spectrum of a reduced complete noetherian local ring $A$. Let $S'$ be the spectrum of the normalization $A'$ of $A$, which is finite over $A$ by Nagata (EGA $0_{\mathrm{IV}}$ 23.1.5); $S'\to S$ is an epimorphism, since $A$ maps injectively into $A'$. Suppose then that there exists a torus $T'$ of $G_{S'}$ with underlying space $E_{S'}$. The two inverse images of $T'$ in $G_{S'\times_S S'}$ are two central subtori having the same underlying space, hence agreeing (Exp. IX 5 bis); so by 3.7, $T'$ comes from a central subtorus $T$ of $G_S$ which obviously has $E$ as underlying space. It therefore suffices to prove the existence of $T'$, which reduces one to the case where $S$ is normal and completes the proof of I).
% typo?: source reference IX 5 bis is retained.

\noindent\textbf{II)} Proof of $\text{(iii)}\implies\text{(ii)}$ when $S$ is normal.

We may restrict to the case where $S$ is integral; let $t$ be its generic point. For every integer $n$ equal to a power of $q$, let ${}_nG$ be the subprescheme of $G$ which is the ``kernel'' of raising to the $n$th power in $G$. Since $E$ is locally closed in $G$ (by (iii) b)), the intersection of $\mathrm{ens}({}_nG)$ with $E$ is locally closed in $\mathrm{ens}(G)$; denote then by $E(n)$ the reduced subprescheme of $G$ with underlying space $\mathrm{ens}({}_nG)\cap E$.

Let us show that the structural morphism $E(n)\to S$ is separated and universally open. For these two properties, one has a valuative criterion (EGA II 7.2.3 and EGA IV 14.5.8). Let therefore $S'$ be an $S$-scheme which is the spectrum of a complete discrete valuation ring with algebraically closed residue field. We have proved that 3.1 (iii) $\implies$ 3.1 (iv), so there exists a subtorus $T'$ of $G_{S'}$ with underlying space $E_{S'}$. Now ${}_nT'$ is finite and \'etale over $S'$, hence is separated and universally open over $S'$, and the same holds for $E(n)\times_S S'$, which has the same underlying space as ${}_nT'$.

Moreover, the fibers of $E(n)$ have the same number of geometric points, namely $r^n$ if $r$ is the rank of the torus $T_t$. Finally, the generic fiber $E(n)_t$, being reduced, is equal to ${}_nT_t$, hence is \'etale over $\Spec(t)$. Since $S$ is normal, it then follows from SGA I 10.11 that $E(n)$ is an \'etale covering of $S$.
% typo?: source gives r^n for the number of geometric points, kept as printed.

If $s$ is a point of $S$, $E(n)_s$ is \'etale over $\Spec\kappa(s)$, hence is reduced and consequently agrees with the group of multiplicative type ${}_nT_s$. Let us show that $E(n)$ is a group subprescheme of $G$. Indeed, let $m$ be the morphism
\[
E(n)\times_S E(n)\to G
\]
induced by multiplication in $G$. The map $\mathrm{ens}(m)$ underlying $m$ factors through $E(n)$, so $m$ factors through the prescheme $E(n)$, since the first member $E(n)\times_S E(n)$ is \'etale over $S$, hence reduced. It then follows from Exp. X 4.8 a) that $E(n)$ is a subgroup of multiplicative type. As already noted (3.2 a)), the family of subgroups $E(n)$ is necessarily coherent. We have therefore proved that (iii) $\implies$ (ii) when $S$ is normal.

\noindent\textbf{III)} Proof of $\text{(ii)}\implies\text{(i)}$.

In fact, we shall show that there exists a unique subtorus $T$ of $G$ with underlying space equal to $E$ and such that ${}_nT=M_n$ for every $n$ equal to a power of $q$. The uniqueness of $T$ follows simply from Exp. IX 4.8 b). To establish the existence of $T$, taking uniqueness into account, we may suppose successively that:

\noindent a) $S$ is noetherian.

\noindent b) The $M_n$ are central subgroups. Indeed, it suffices to replace $G$ by $Z_n=\uCentr_G(M_n)$ for $n$ sufficiently large (2.5 and 2.5 bis).

\noindent c) $S$ is the spectrum of a local ring. Suppose indeed that the problem is solved after every base change $\Spec(\OO_{S,s})\to S$, where $s$ ranges over the points of $S$. Let $T_s$ be the subtorus of $G\times_S\Spec(\OO_{S,s})$ thus obtained. For every $s$, there then exist an open neighborhood $U$ of $s$ and a subtorus $T$ of $G|_U$ extending $T_s$ (Exp. VIB \S~10 and 3.6). We have proved that 3.1 (ii) $\implies$ 3.1 (iv); since $S$ is noetherian, $E$ is therefore constructible (3.3). Consequently, after restricting $U$, one may suppose that $\mathrm{ens}(T)=E\times_S U$ (EGA IV 9.5.2). But then, for every integer $n$ equal to a power of $q$, ${}_nT$ and $M_n\times_S U$ are two subgroups of multiplicative type of $G|_U$ having the same fibers, hence agreeing, since $M_n$ is central (Exp. IX 5.3 bis). In brief, under the hypotheses made, there exists a solution locally on $S$, hence, because of uniqueness, there exists a global solution.

\noindent d) $S$ is the spectrum of a complete noetherian local ring and, if $s$ is the closed point, $T_s$ is trivial. This follows from EGA $0_{\mathrm{III}}$ 10.3.1 and fpqc descent.

\noindent e) $S$ is reduced. One applies 2.2.

\noindent f) $S$ is normal. One applies Nagata's finiteness theorem (EGA $0_{\mathrm{IV}}$ 23.1.5) and Lemma 3.7 in the case of central subtori.

These reductions having been made, since $T_s$ is trivial, $(M_n)_s$ is trivial, hence $M_n$ is trivial (Exp. X 3.3). If $t$ is a point of $S$, the subgroups ${}_nT_t$ of $T_t$ are therefore trivial for every $n$ equal to a power of $q$, and it follows easily from Exp. X 1.4 that $T_t$ itself is trivial. It then suffices to prove the lemma:
\begin{lemma}{3.10}\label{XV.3.10}
Under the hypotheses of 3.1 (ii), suppose in addition that $S$ is noetherian and normal, and that for every point $s$ of $S$, $T_s$ is a trivial torus. Then there exists a unique subtorus $T$ of $G$ with underlying space equal to $E$, such that ${}_nT=M_n$ for every $n$ equal to a power of $q$. Moreover, $T$ is trivial.
\end{lemma}
\begin{proof}
The uniqueness of $T$ follows from the fact that $T$, being smooth over $S$, is reduced. To prove existence, one may suppose $S$ irreducible with generic point $\eta$. Let $r$ be the rank of $T_\eta$, $T'=\mathbf G_{\mathrm m,S}^r$, $u_\eta$ an isomorphism of $T'_\eta$ onto $T_\eta$, ${}_nu_\eta$ the restriction of $u_\eta$ to ${}_nT'_\eta$. Since ${}_nT'$ and $M_n$ are trivial, there exists a unique extension ${}_nu$ of ${}_nu_\eta$ to an $S$-isomorphism of ${}_nT'$ onto $M_n$. I say that for every point $s$ of $S$, there exists a group isomorphism, necessarily unique:
\[
u_s:T'_s\isomto T_s
\]
extending ${}_nu_s$ for every $n$ equal to a power of $q$. Indeed, let $S_1$ be an $S$-scheme, the spectrum of a complete discrete valuation ring with algebraically closed residue field, whose generic point $t_1$ projects to $\eta$ and whose closed point $s_1$ projects to $s$ (EGA II 7.1.9). It follows from the proof of 3.1 (ii) $\implies$ 3.1 (iv) that there exists a subtorus $T^1$ of $G_{S_1}$ such that $(M_n)_{S_1}={}_nT^1$ for every $n$ equal to a power of $q$. Since $S_1$ is normal, $T^1$ is isotrivial (Exp. X 5.16), and it follows from the classification of isotrivial tori (Exp. X 1.2) and SGA1 V 8.2 that $T^1$, having trivial generic fiber, is trivial. One can therefore extend the isomorphism $u_\eta\times_\eta t_1$ to an $S_1$-isomorphism of group schemes:
\[
u^1:T'\times_S S_1\isomto T^1.
\]
The restriction of $u^1$ to ${}_nT'\times_S S_1$, on the one hand, and ${}_nu\times_S S^1$, on the other hand, agree on the generic fiber by construction, hence agree. The restriction $u^1_{s_1}$ of $u^1$ to the closed fiber therefore gives the required extension after extension of the residue field $\kappa(s)\to\kappa(s_1)$. It then follows from Exp. IX 4.8 a) and fpqc descent that $u^1_{s_1}$ descends to $\kappa(s)$ as a group isomorphism $u_s:T'_s\isomto T_s$ answering the question.
% typo?: source writes S^1 here although the scheme was denoted S_1.

Moreover, $T'$ is smooth over $S$, which is normal, hence is normal. It then follows from an easy criterion for extending rational maps (EGA IV$_4$ 20.4.6) that, for there to exist an $S$-morphism $u:T'\to G$ whose restriction to $T'_s$, for every point $s$ of $S$, is the composite morphism:
\[
T'_s\xrightarrow{u_s}T_s\to G_s,
\]
it is necessary and sufficient that this be so after every base change $S_1\to S$, where $S_1$ is the spectrum of a complete discrete valuation ring with algebraically closed residue field.

Now in the present case, an argument analogous to the one just made shows that the extension condition is indeed satisfied. Denote by $u$ the $S$-morphism $T'\to G$ thus obtained.

Let us show that $u$ is indeed a group morphism. Let $m_{T'}$ (resp. $m_G$) be the morphism defining multiplication in $T'$ (resp. $G$). One must verify that $m_G\circ(u\times_S u)=u\circ m_{T'}$. Now the coincidence subprescheme of these two morphisms is a subprescheme of $T'\times_S T'$ containing the fibers (since $u_S$ is a group morphism), hence having the same underlying space as $T'\times_S T'$, hence equal to the latter, since $T'\times_S T'$ is smooth over $S$, thus reduced.
% typo?: source says u_S rather than u_s, kept as printed.

Finally, note that $u$ is a monomorphism (since this is so on the fibers), hence is an immersion (Exp. VIII 7.9). The image of $T'$ under $u$ is then a subtorus of $G$ having $E$ as underlying set.
\end{proof}
This completes the proof of Theorem 3.1.

% original p. 398
\sgasection{4}{\texorpdfstring{Characterization of a subtorus $T$ by the subgroups ${}_nT$}{Characterization of a subtorus T by the subgroups nT}}
\par\medskip\noindent\textit{1. Statement of the main theorem}\par
\begin{theorem}{4.1}\label{XV.4.1}
Let $S$ be a connected locally noetherian prescheme, $G$ a group $S$-prescheme of finite type over $S$, $q$ an integer $>1$ invertible on $S$, $r$ a positive integer. For every integer $n$ equal to a power of $q$, let $M(n)$ be a group subscheme of $G$, of multiplicative type and of type $(\ZZ/n\ZZ)^r$. One supposes that:

\noindent a) The family of subgroups $M(n)$ is coherent, that is, if the integer $m$ divides $n$, one has:
\[
{}_mM(n)=M(m).
\]
\noindent b) There exist a point $s$ of $S$ and a subtorus $T_s$ of $G_s$ such that:
\[
M(n)_s={}_nT_s\quad\text{for every }n.
\]
\noindent c) For every point $t$ of $S$, there exists a closed affine subscheme $F_{\bar t}$ of $G_{\bar t}$ containing $M(n)_{\bar t}$ for every $n$.

Then there exists one and only one subtorus $T$ of $G$ such that ${}_nT=M(n)$ for every $n$ equal to a power of $q$.
\end{theorem}

One has an analogous theorem concerning the lifting of morphisms:
\begin{theorem}{4.1 bis}\label{XV.4.1bis}
Let $S,G,q$ be as above, $T$ an $S$-torus, and for every integer $n$ equal to a power of $q$, let $u(n)$ be a group $S$-morphism ${}_nT\to G$. One supposes that:

\noindent a) The family of morphisms $u(n)$ is coherent, i.e. if $m$ divides $n$, one has:
\[
u(m)=u(n)|_{{}_mT}.
\]
\noindent b) There exist a point $s$ of $S$ and a group morphism:
\[
u_s:T_s\to G_s
\]
such that $u_s|_{{}_nT_s}=u(n)_s$ for every $n$ equal to a power of $q$.

\noindent c) For every point $t$ of $S$, there exists a closed affine subscheme $F_{\bar t}$ of $G_{\bar t}$ containing $u(n)_{\bar t}({}_nT_{\bar t})$ for every $n$.

Then there exists a unique group morphism $u:T\to G$ such that for every $n$ equal to a power of $q$, the restriction of $u$ to ${}_nT$ is equal to $u(n)$.
\end{theorem}
\begin{remark}{4.2}\label{XV.4.2}
Using the lower semicontinuity of the abelian rank of a group prescheme flat of finite type over the spectrum of a discrete valuation ring (cf. Exp. X 8.7), one may, in the statements of 4.1 and 4.1 bis, weaken condition c) by merely asking that the required closed affine subscheme $F_{\bar t}$ exist for every maximal point $t$ of $S$.
\end{remark}

\par\noindent\emph{Let us show how 4.1 bis follows from 4.1.} Let $G'=G\times_S T$. For every integer $n$ equal to a power of $q$, consider the group morphism:
\[
v(n):{}_nT\to G'
\]
whose projections onto $G$ and $T$ are respectively $u(n)$ and the canonical immersion ${}_nT\to T$. The morphism $v(n)$ is therefore an immersion; let $M(n)$ be the image subgroup.
% typo: respectivment -> respectively.
It is clear that the family of subgroups $M(n)$ is coherent in the sense of 4.1, that the group $M(n)_s$ is equal to ${}_nT'_s$, where $T'_s$ is the subtorus of $G'_s$ which is the graph of $u_s$, and that for every point $t$ of $S$, the closed affine subscheme $F_{\bar t}\times_{\bar t}T_{\bar t}$ of $G'_{\bar t}$ contains the subgroups $M(n)_{\bar t}$ for every $n$. By 4.1, there therefore exists a subtorus $T'$ of $G'$ such that ${}_nT'=M(n)$ for every $n$ equal to a power of $q$. Let $f$ be the restriction to $T'$ of the projection of $G'$ onto $T$, and $f(n)$ the restriction of $f$ to $M(n)$. One has:
\[
f(n)\circ v(n)=\mathrm{id}_{{}_nT}.
\]
The fiber at $s$ of $T'$ is the torus already denoted by $T'_s$, equal to the graph of $u_s$ (this follows from Exp. IX 4.8 b)), so $f_s$ is an isomorphism. But $\Ker f$ and $\Coker f$ are groups of multiplicative type (Exp. IX 2.7), of constant type since $S$ is connected, hence reduced to the unit group, and $f$ is an isomorphism. Let $v$ be the inverse isomorphism of $f$. One has:
\[
v|_{{}_nT}=v(n).
\]
Consequently, the composite of $v$ and the projection of $G'$ onto $G$ is a morphism $u:T\to G$ answering the question. What precedes proves the existence of the morphism $u$; as for uniqueness, it follows in any case from Exp. IX 4.8 a).

\par\medskip\noindent\textit{2. Application}\par
We propose to generalize Theorem 7.1 of Exp. IX.

Let $A$ be a complete noetherian local ring, $I$ its maximal ideal, $S=\Spec A$, $S_m=\Spec(A/I^m)$. For every prescheme $X$, put $X_m=X\times_S S_m$.

Let then $G$ be a group $S$-prescheme of finite type, $T$ an $S$-torus, $q$ an integer invertible on $S$, and $u_m:T_m\to G_m$ ($m\in\NN$) a coherent family of group morphisms. With the integer $n$ ranging over the powers of $q$, denote by $u_m(n)$ the restriction of $u_m$ to ${}_nT_m$, and by $u(n)$ the unique group morphism:
\[
u(n):{}_nT\to G
\]
extending the morphisms $u_m(n)$ for every $m$ (1.6 a)). We shall say that the family $(u_m)$, $m\in\NN$, is admissible if for every point $t$ of $S$, there exists a closed affine subprescheme $F_{\bar t}$ of $G_{\bar t}$ containing $u(n)_{\bar t}({}_nT_{\bar t})$ for every $n$ equal to a power of $q$ (this property is independent of $q$, as we shall see).
\begin{proposition}{4.3}\label{XV.4.3}
With the notation above, the canonical map:
\[
\Hom_{S\text{-gr}}(T,G)\to\varprojlim_m\Hom_{S_m\text{-gr}}(T_m,G_m)
\]
induces an isomorphism from the first member onto the subset of the second member formed by the ``admissible'' coherent families.
\end{proposition}
Indeed, it suffices to apply 4.1 bis, taking as $s$ the closed point of $S$.
\begin{corollary}{4.4}\label{XV.4.4}
With the notation above, suppose in addition that $G$ has affine fibers; then the canonical map:
\[
\Hom_{S\text{-gr}}(T,G)\to\varprojlim_m\Hom_{S_m\text{-gr}}(T_m,G_m)
\]
is an isomorphism.
\end{corollary}
Indeed, if $G$ has affine fibers, every coherent family is ``admissible''.
\begin{remark}{4.5}\label{XV.4.5}
When $G$ is separated, in 4.3 one may replace the ring $A$ by a noetherian $I$-adic ring; indeed, one may use EGA III 5.4.1 instead of 1.6 a).
\end{remark}

\par\medskip\noindent\textit{3. Proof of 4.1}\par
\begin{lemma}{4.6}\label{XV.4.6}
Let $k$ be a field, $G$ an algebraic $k$-group, $q$ an integer prime to the residual characteristic of $k$, $r$ an integer $>0$, $M(n)$ (with $n$ ranging over the powers of $q$) a coherent family of subgroups of $G$, of multiplicative type and of type $(\ZZ/n\ZZ)^r$. Then there exists a smallest closed subscheme $H$ of $G$ containing the $M(n)$ for every $n$. Moreover, $H$ is an algebraic subgroup of $G$, smooth, connected and commutative, ``whose formation is compatible with extension of the base field''.
\end{lemma}
Let $M$ be the subset of $\mathrm{ens}(G)$ which is the union of the underlying sets of the subgroups $M(n)$ for every $n$, and $H$ the reduced closed subscheme of $G$ whose underlying space is the closure of $M$. Since $M(n)$ is \'etale over $k$, hence reduced, $M(n)$ is contained in $H$, and consequently $H$ is the smallest closed subscheme of $G$ containing $M(n)$ for every $n$. Let now $\bar k$ be an algebraically closed extension of $k$. By construction, the subschemes $M(n)$ are schematically dense in $H$ (Exp. IX 4.1), so the $M(n)_{\bar k}$ are schematically dense in $H_{\bar k}$ (Exp. IX 4.5); moreover, $M(n)_{\bar k}$ is reduced, and it follows that $H_{\bar k}$ is necessarily equal to the closed reduced subscheme of $G_{\bar k}$ whose underlying space is the closure of $M_{\bar k}$. This proves that $H$ is geometrically reduced (EGA IV 4.6.1). Since the family $M(n)$ is coherent, $M$ is stable under the group law; moreover, $M\times_k M$ is dense in $\mathrm{ens}(H\times_k H)$ and $H\times_k H$ is reduced (EGA IV 4.6.5); one immediately deduces that $H$ is an algebraic subgroup of $G$. Moreover, $H$ is smooth over $S$, since it is geometrically reduced (Exp. VIA 1.3.1), and $H$ is commutative since the $M(n)$ are commutative. It remains to see that $H$ is connected.
% typo?: source says smooth over S rather than k, kept as printed.
